\documentclass{article}
\usepackage{graphicx} % Required for inserting images
\usepackage{caption} % Required for \captionof
\usepackage{amsmath}
\usepackage{float}
\usepackage{comment}
\usepackage[T1]{fontenc}
\usepackage[utf8]{inputenc}
\usepackage{lmodern}
\usepackage{microtype}
\usepackage{xcolor}
\usepackage{framed}
\renewenvironment{leftbar}{%
  \def\FrameCommand{\textcolor{orange}{\vrule width 3pt} \hspace{10pt}}%
  \MakeFramed {\advance\hsize-\width \FrameRestore}}%
 {\endMakeFramed}


\begin{document}

\section*{7. Feladat: Az alumínium fajhőjének mérése}

Ebben a feladatban az alumínium fajhőjét határoztuk meg. Ehhez először is megmértük a kapott alumínium henger paramétereit, hogy térfogatának kiszámolásával megbecsülhessük a tömegét. A henger átmérője 40 mm-nek, a magassága pedig 150 mm-nek adódott. Ebből a számított térfogat 189 cm$^3$-nek adódott. Az alumínium sűrűségét 2,7 g/cm$^3$-nek véve a becsült tömeg 509 g lett. A tömeget lemértük mérleggel is, és így 534 g-ot kaptunk. Ez viszonylag nagyobb eltérés. A továbbiakban a mérleggel mért tömeget használtuk a számolásokban a pontosabb eredményekért. 

Ezek után mindkét termoszt feltöltöttük annyi vízzel (az egyiket hideggel, a másikat meleggel), hogy ha az alumínium hengert belelógatjuk, akkor teljesen elmerüljön benne, de a víz ne csorduljon ki (itt már a piros termoszt használtuk melegvizesként). A folyadékok tömegét megmértük, majd a hideg vízbe beletettük az alumínium hengert. Megvártuk, amíg beáll a termikus egyensúly, majd leolvastuk a kezdeti hőmérsékleteket. A kezdeti adatok a következőknek adódtak: a hideg víz tömege 415 g, hőmérséklete 25,45 $^{\circ}$C. A meleg termoszban lévő víz tömege szintén 415 g volt, kezdeti hőmérséklete 92,2 $^{\circ}$C. Ezek után a lehűlt hengert áttettük a melegvizes termoszba, és addig mértük a hőmérsékletadatokat, amíg az új termikus egyensúly be nem állt. Ezt az alábbi grafikonon ábrázoltuk. Az új egyensúlyi hőmérsékletet 75,8 $^{\circ}$C-nak mértük. A fent megadott adatokból és a végső egyensúlyi hőmérsékletből pedig kiszámoltuk az alumínium fajhőjét, a leírásban mutatott módszerrel. Megint feltettük azt a számolások során, hogy a környezettel nincsen hőcsere, és megint a víz irodalmi fajhőjét használtuk (4178 J/(kg$^{\circ}$C)). A piros termosz hőkapacitását a közelített 95 J/K-nek vettük. 

% ! szamitas leirasa

% $d_{\text{henger}} = 40,0 \text{ mm}$\\
% $h_{\text{henger}} = 15 \text{ cm}$\\
% $V_{\text{henger}} = 188,5 \text{ cm}^3$\\
% $\rho_{\text{Al}} = 2,7 \text{ g/cm}^3$\\
% $m_{\text{becsült}} = 509 \text{ g}$\\
% $m_{\text{mért}} = 534 \text{ g}$\\
% A hideg termosz adatai:\\
% $m_{\text{víz}} = 415 \text{ g}$\\
% $T_{\text{hideg}} = 25,45^{\circ}\text{C}$\\
% A meleg termosz adatai:\\
% $m_{\text{víz}} = 415 \text{ g}$\\
% $T_{\text{víz}} = 92,2^{\circ}\text{C}$\\
% $T_{\text{egyensúlyi}} = 75,8^{\circ}\text{C}$\\
% $c_{\text{víz}} = 4178 \text{ J/(kg}^{\circ}\text{C)}$ \\  % irodalmi adat
% $C_{\text{termosz}} = 95 \text{ J/K}$\\
% $c_{\text{Al}} = 634 \text{ J/(kg}^{\circ}\text{C)}$\\

\begin{figure}[H]
    \centering
    \includegraphics[width=0.8\textwidth]{7_1_graf.pdf}
    \caption*{7. feladat: Mért hőmérséklet az idő függvényében}
    \label{fig:7_1}
\end{figure}

Az eredményünk ugyan nem lett tökéletes, de nem is volt annyira távol az irodalmi adattól. A mi mérésünk (és számolásunk) alapján az alumínium fajhője 634 J/(kg$^{\circ}$C)-nak adódott. Ez elég jelentős mérési hibával és bizonytalansággal terhelt, így 630 J/(kg$^{\circ}$C)-ra kerekítettünk. Az irodalmi adat 890 J/(kg$^{\circ}$C). Ez 29 $\%$-os hibát jelent. A mérés eredménye a körülményekhez képest elfogadható. A legnagyobb hibaforrás valószínűleg az volt, hogy a második alkalommal kapott piros termosz jóval kisebb volt, mint a feketék vagy a zöldek, így a hőmérő egészen az alumínium henger falához nyomódott, ez a homersekletmerest reszben befolyasolhatta. Emellett a piros termosz kis mérete miatt a keverés is nehézkesebben történt meg, így jóval pontatlanabb volt a hőmérséklet-mérés is. Nyilván valamennyi környezet felé történő hőszivárgás is fellépett, de ez nem volt számottevő: az egyensúlyi hőmérséklet beállta után a hőmérséklet már nem csökkent tovább, igazolva ezzel, hogy a mérés ideje alatt a hőveszteség kicsi volt (ez egybevág a 3. feladat eredményeivel is). A termosz hőkapacitásának becsült értéke viszont nagyon nagy hibával terhelt, így a relatív nagy eltérés egy jelentős része valószínűleg ebből a pontatlanságból adódott. A mérés hossza miatt itt is hibaforrásként szerepet játszik a hő környezetbe való távozása. Ez a harmadik feladat eredményei alapján a mérés hosszát tekintve akár a Celsius nagyságrandbe is eshet, ezzel jelentős hibát adva a mérésünk eredményének.
% !homersekleti korrekcio

Hibaszámítás: .........

\section*{8. Feladat: A sárgaréz fajhőjének mérése}

Ebben a feladatban az előzőhöz nagyon hasonló mérést végeztünk, azzal a különbséggel, hogy most nem alumínium, hanem sárgaréz fajhőjét állapítottuk meg. A réz henger paramétereit ugyanakkorának mértük, mint az alumínium hengerét. A továbbiakban is megegyezett a mérés menete az előző feladatéval. A mért és számított adatainkat az alábbiakban ismertetjük:

\begin{itemize}
    \item \textbf{A rézhenger paraméterei:}
    \begin{itemize}
        \item Átmérő: $d_{\text{henger}} = 40 \text{ mm}$
        \item Magasság: $h_{\text{henger}} = 15 \text{ cm}$
        \item Sűrűség: $\rho_{\text{henger}} = 8,73 \text{ g/cm}^3$
        \item Becsült tömeg: $m_{\text{becsült}} = 1646 \text{ g}$
        \item Mért tömeg: $m_{\text{mért}} = 1592 \text{ g}$
    \end{itemize}
    \item \textbf{A hideg termosz adatai:}
    \begin{itemize}
        \item Kezdeti hőmérséklet: $T_{\text{hideg}} = 26,2^{\circ}\text{C}$
    \end{itemize}
    \item \textbf{A meleg termosz adatai:}
    \begin{itemize}
        \item Víz tömege: $m_{\text{víz}} = 436 \text{ g}$
        \item Kezdeti hőmérséklet: $T_{\text{víz}} = 82,3^{\circ}\text{C}$
        \item Egyensúlyi hőmérséklet: $T_{\text{egyensúlyi}} = 66,5^{\circ}\text{C}$
    \end{itemize}
    \item \textbf{Fajhő és hőkapacitás adatok:}
    \begin{itemize}
        \item Víz fajhője (irodalmi): $c_{\text{víz}} = 4178 \text{ J/(kg}^{\circ}\text{C)}$
        \item Rozsdamentes acél fajhője: $c_{\text{rozsdamentes acél}} = 500 \text{ J/(kg}^{\circ}\text{C)}$
        \item Termosz hőkapacitása: $C_{\text{termosz}} = 95 \text{ J/K}$
        \item \textbf{Sárgaréz kiszámolt fajhője: $c_{\text{Cu}} = 472 \text{ J/(kg}^{\circ}\text{C)}$}
    \end{itemize}
\end{itemize}

\begin{figure}[H]
    \centering
    \includegraphics[width=0.8\textwidth]{8_1_graf.pdf}
    \caption*{8. feladat: Mért hőmérséklet az idő függvényében}
    \label{fig:8_1}
\end{figure}

A sárgaréz fajhője 472 J/(kg$^{\circ}$C)-nak adódott, amelyet a bizonytalanságok tükrében kerekítve, 470 J/(kg$^{\circ}$C)-ként érdemes figyelembe venni. Ez a 385 J/(kg$^{\circ}$C)-os irodalmi adattal viszonylag jó egyezésben van, a fajhőt kb. 22\%-os relatív hibával határoztuk meg. Bár ez az eredmény ezúttal sem tökéletes, a körülményeket tekintve nem is nagyon távoli. A hibák lehetséges okai hasonlóak az előző mérésnél tapasztaltakhoz. A rézhengernek már fix felfüggesztése volt (nem kötél, hanem merev drót), így amikor belehelyeztük a melegvizes termoszba, még jobban akadályozta a mágneses keverő mozgását. Emellett a hőmérő továbbra is a henger oldalához szorult, így nem a folyadék átlagos hőmérsékletét mértük. Ezeken felül itt is nagyon jelentős hibaforrást jelent a termosz hőkapacitásának meglehetősen nagy bizonytalansága. A mérés hossza miatt itt is hibaforrásként szerepet játszik a hő környezetbe való távozása. Ez a harmadik feladat eredményei alapján a mérés hosszát tekintve akár a Celsius nagyságrandbe is eshet, ezzel jelentős hibát adva a mérésünk eredményének.

Hibaszámítás: .........

\section*{9. Feladat: A jég olvadáshőjének mérése keveréssel}

Ebben a feladatban a jég olvadáshőjét próbáltuk megállapítani olyan módon, hogy a piros termoszunkba meleg vizet töltöttünk, és megvártuk, míg ebben a termoszban beáll a termikus egyensúly. Ekkor körülbelül 0 $^{\circ}$C-os hőmérsékletű jeget tettünk a meleg vízbe, és elkezdtük gyűjteni az adatokat. Ezeket később grafikon formájában ábrázoltuk. A mérést addig folytattuk, míg a termoszban lévő összes jég el nem olvadt. A mérés végén a termoszban lévő összes víz tömegének megmérésével meghatároztuk a jég tömegét. Ezekből az adatokból pedig kiszámoltuk a jég olvadáshőjét. A számolásokhoz a víz fajhőjére az irodalmi adatot használtuk, és feltettük, hogy a rendszer és a környezet közti hőcsere elhanyagolható. A mért és számolt eredményeink alább láthatóak.

\begin{itemize}
    \item \textbf{Mért tömegek:}
    \begin{itemize}
        \item Víz tömege: $m_{\text{víz}} = 477 \text{ g}$
        \item Teljes tömeg (jég elolvadása után): $m_{\text{teljes}} = 510 \text{ g}$
        \item Jég tömege: $m_{\text{jég}} = 33 \text{ g}$
    \end{itemize}
    \item \textbf{Hőmérsékleti adatok:}
    \begin{itemize}
        \item Kezdeti hőmérséklet: $T_{\text{kezdeti}} = 93,2^{\circ}\text{C}$
        \item Egyensúlyi hőmérséklet: $T_{\text{egyensúlyi}} = 80,8^{\circ}\text{C}$
    \end{itemize}
    \item \textbf{Fajhő és állandók:}
    \begin{itemize}
        \item Víz fajhője (irodalmi): $c_{\text{víz}} = 4178 \text{ J/(kg}^{\circ}\text{C)}$
        \item Termosz hőkapacitása: $C_{\text{termosz}} = 95 \text{ J/K}$
        \item \textbf{A jég kiszámolt olvadáshője: $L_{f} = 450 \text{ kJ/kg}$}
    \end{itemize}
\end{itemize}

\begin{figure}[H]
    \centering
    \includegraphics[width=0.8\textwidth]{9_1_graf.pdf}
    \caption*{9. feladat: Jég olvadáshő}
    \label{fig:9_1}
\end{figure}

A hiba leginkább abból származhatott, hogy a jég hőmérsékletét nem tudtuk pontosan, és nem kizárt, hogy hidegebb volt, mint 0 $^{\circ}$C, így már az olvadás előtt is hőt vont el a víztől. Emellett természetesen megint szerepet játszott a termosz hőkapacitásának pontatlansága. És valamennyire a disszipálódhatott is hő, bár ez inkább a tized Celsius nagyságrandbe esik (a 3. feladat eredményei alapján).

Hibaszámítás: .........


\end{document}